% RESEARCH_DRAFT_SHA256: 65b2263d2a57f5665035eccdafeea51de3537e117f4e977f1e5edb6007aa449b
% Layout and source-reference wrappers expanded for the web renderer.
\documentclass[11pt]{article}
\usepackage[margin=25mm]{geometry}
\usepackage{amsthm}
\usepackage{fontspec,unicode-math}
\setmainfont{Latin Modern Roman}
\setmathfont{Latin Modern Math}
\usepackage{amsmath,amssymb,mathtools}
\usepackage{xurl,hyperref}
\hypersetup{colorlinks=true,urlcolor=blue,linkcolor=blue}
\setcounter{secnumdepth}{0}
\setlength{\parindent}{0pt}
\setlength{\parskip}{5pt}
\setlength{\emergencystretch}{3em}
\newcommand{\sref}[2]{\hyperlink{source:#1:#2}{[S#2]}}
\newcommand{\cataloguescope}[1]{\paragraph{Source question.}#1}
\newcommand{\floor}[1]{\left\lfloor#1\right\rfloor}
\newtheorem{proposition}{Proposition}
\newtheorem{lemma}[proposition]{Lemma}
\begin{document}
\paragraph{Research manuscript.} OpenAI \texttt{gpt-6-astra}, reasoning \texttt{ultra}; 11 October 2026.
\paragraph{Preserved source body.} The complete catalogue mathematical text and sources follow with layout and source links expanded, including source qualifications and its historical preparation notes. New research begins at the section headed Research attempt.

\section{Benign, Tempered, and Catastrophic Overfitting Regimes}\label{statement}

\subsection{Definitions and mathematical statement}

\paragraph{Shared notation.}Standard number systems are denoted by $\mathbb N,\mathbb Z,\mathbb Q,\mathbb R,\mathbb C$; index ranges, computational models and specialized notation follow the cited problem source and the definitions below.


An interpolating estimator fits the training labels exactly. Its prediction risk is the expectation of a specified loss on a fresh data point, and excess risk is the difference from the source population benchmark. In the source terminology, benign overfitting has vanishing asymptotic excess risk; tempered overfitting has a positive finite limiting excess risk; catastrophic overfitting has divergent excess risk. A kernel method uses a positive-semidefinite kernel $K(x,x')$ and its associated function space; neural-network interpolators instead use the specified network class and training rule.
\paragraph{Statement recorded in the supplied source.}
Characterize when interpolation and overfitting in neural networks or kernel methods is benign, tempered, or catastrophic. \href{https://doi.org/10.4171/owr/2023/41}{[S2]}
\paragraph{Scope qualification.}
These regimes refer to the source risk and asymptotic scaling. Training interpolation alone does not specify a regime, and different loss or dimension limits are not identified.
\paragraph{Source status.}
The supplied catalogue labels this record partially solved. Its source-status metadata was checked against the supplied snapshot on 2026-10-11. This import does not independently establish current open status or mathematical correctness. The catalogue literature assessment is dated 2026-08-22.
\subsection{Short English statement}
Determine which data, kernels and neural-network learning rules lead to vanishing, bounded positive, or divergent errors after interpolation.
\subsection{Sources}
\begin{itemize}
\item[S1] Ulam AI and the UnsolvedMath Contributors, UnsolvedMath: A Curated Collection of Open Mathematics Problems, catalogue entry. Catalogue curation and metadata: CC BY 4.0. The supplied source question is retained; definitions and the English summary are editorial additions.
\url{https://huggingface.co/datasets/ulamai/UnsolvedMath}.
\item[S2] Mini-Workshop: Interpolation and Over-parameterization in Statistics and Machine Learning​ (2024), Oberwolfach Reports, DOI 10.4171/owr/2023/41, pp. 2359--2376
\url{https://doi.org/10.4171/owr/2023/41}.
\item[S3] Mallinar et al., overfitting taxonomy
\url{https://papers.neurips.cc/paper_files/paper/2022/hash/08342dc6ab69f23167b4123086ad4d38-Abstract-Conference.html}.
\item[S4] benign overfitting in two-layer networks
\url{https://proceedings.neurips.cc/paper_files/paper/2022/file/a12c999be280372b157294e72a4bbc8b-Paper-Conference.pdf}.
\textit{Underlying Oberwolfach Reports material retains its source-specific licence, including CC BY-SA 4.0 where applicable; see THIRD\_PARTY\_NOTICES.md.}
\end{itemize}
\paragraph{Preparation and review.}
Original question attributed to the cited source and UnsolvedMath contributors. Codex (GPT-6; exact serving revision not exposed) assisted with definitions, English summaries, source inspection and formatting. No new human mathematical authorship or independent proof verification is asserted. Catalogue identity checked against the complete supplied registry and existing repository statements; known flagged repetitions and obvious identical conjectures excluded. Full mathematical-equivalence/literature review has not been performed. Statement preparation and final review followed contribution prompts C and D; automated checks verify structure and identity, not open status.
\label{end:statement}

\section{Research attempt}
\subsection{Literature, source context and scope}
The full original catalogue body and source qualifications are preserved above. The relevant primary workshop passages were read in the cached publisher report. Historical claims and subsequent attributed results are distinguished from the mathematical work below. A restricted example, standard lemma or known theorem is not counted as new progress on the entire original target. Missing assumptions and repeated or rhetorical extracts are not silently turned into different solved problems.

\subsection{Three exact averaged-risk regimes in a linear kernel model}
The source and the Mallinar et al. taxonomy distinguish several asymptotic behaviors. We specify which risk is computed: expectation over training design and noise as well as the fresh test point. This need not agree with an in-probability classification when tails are heavy. Let $X\in\mathbb R^{n\times p}$ have independent standard Gaussian entries, $p>n+1$, and let $y=\xi$ be independent mean-zero noise with covariance $\sigma^2I_n$, where $\sigma>0$. The target regression function is zero. Use the minimum Euclidean norm interpolator, equivalently interpolation with the linear kernel $K(x,z)=x^Tz$.
\begin{proposition}
Its expected excess squared prediction risk is exactly
\[
\mathbb E R(\widehat\beta)=\frac{\sigma^2 n}{p-n-1}.
\]
\end{proposition}
\begin{proof}
Gaussian rows are linearly independent almost surely: conditional on preceding rows, a new row has a density and a proper linear subspace has Lebesgue measure zero. Thus $\widehat\beta=X^T(XX^T)^{-1}\xi$. A fresh standard Gaussian covariate has covariance identity, so the conditional test excess risk is $\|\widehat\beta\|^2$. Averaging over the noise gives $\sigma^2\operatorname{tr}((XX^T)^{-1})$.
For the first diagonal inverse entry, separate the first row $x_1$ from the remaining matrix $Y$. Solving the block equations for the first inverse column, or eliminating its lower block, gives
\[
[(XX^T)^{-1}]_{11}=
\big[x_1(I-Y^T(YY^T)^{-1}Y)x_1^T\big]^{-1}.
\]
For $n=1$ take the projector in brackets to be $I$. The bracketed matrix is orthogonal projection onto a space of dimension $k=p-n+1$. Conditional on $Y$, rotational invariance of the independent Gaussian row makes the denominator a sum of $k$ independent squared standard Gaussians. Its density is proportional to $t^{k/2-1}e^{-t/2}$ for $t>0$, by polar coordinates in $\mathbb R^k$. Its reciprocal mean is the ratio of the integrals with exponents $k/2-2$ and $k/2-1$. Integration by parts gives
\[
\int_0^\infty t^{k/2-1}e^{-t/2}\,dt
=(k-2)\int_0^\infty t^{k/2-2}e^{-t/2}\,dt,
\]
where the boundary terms vanish since $k>2$. Thus each diagonal inverse entry has mean $1/(p-n-1)$. Summing them proves the formula.
\end{proof}
For $n\to\infty$, the choices $p=n^2+n+1$, $p=2n+1$, and $p=n+2$ respectively give averaged excess risks $\sigma^2/n$, $\sigma^2$, and $\sigma^2 n$. Each method exactly interpolates noisy labels almost surely. Therefore interpolation alone cannot choose among even these three expected-risk regimes. At or below $p=n+1$, the displayed finite expectation calculation is not valid; the inverse moment diverges at the relevant Gaussian boundary.
\subsection{Verification and remaining gaps}
The integral recurrence was checked for small integer degrees of freedom using elementary gamma recurrences, and the three dimension sequences give the stated risks exactly. Numerical full-rank Gaussian designs also verify the conditional trace identity; finite simulations are not used to establish asymptotic expectations. This is a classical solvable linear-kernel model, not a classification for nonlinear networks, arbitrary kernels, optimization rules, or high-probability risk. The missing extension is control of the corresponding data-dependent interpolation operator under those models.

\section{References and source attribution}
Primary workshop report: \url{https://doi.org/10.4171/owr/2023/41}. Relevant passages read: 2023/41 pp. 2365--2366, interpolation questions and dimension-dependent tempered/benign examples.

\url{https://papers.neurips.cc/paper_files/paper/2022/hash/08342dc6ab69f23167b4123086ad4d38-Abstract-Conference.html}. Read the primary taxonomy abstract: empirical three-part taxonomy and kernel-regression specialization; its full kernel theorem is not reproduced.

All retrievals were checked on 11 October 2026. Original question attribution remains in the preserved Sources section. External theorems are attributed and not claimed as new discoveries or fully verified merely by citation.


\paragraph{Integration review.} Codex (GPT-6; exact serving revision not exposed) reviewed the written arguments and expanded layout and source-reference wrappers. Repository metadata and the production web renderer were checked using Python 3.13 and JavaScript with MathJax 3.2.2. This model-assisted review is not independent human proof verification; no PDF compilation or formal proof checking is asserted.

\section{Completion Estimate}
\noindent\textbf{COMPLETION: 0\%}
No substantiated new progress toward the entire original target is claimed. The worked results and failed route are delimited above; a known result, illustrative model, or auxiliary obstruction is not counted as original-target completion.

\section{Final status and attribution}
\textbf{Attempt status: unresolved.} This describes this research attempt, not a new ruling on the source catalogue's status or on external literature claims.
The mathematics and this self-audit were generated by OpenAI \texttt{gpt-6-astra}, reasoning \texttt{ultra}, on 11 October 2026. Codex (GPT-6; exact serving revision not exposed) orchestrated the batch, read and reviewed its written mathematical arguments, and contributed formatting/repository review. No human mathematical authorship, independent human proof checking, or formal verification is claimed. The historical preparation paragraphs in the reproduced catalogue body concern that body only.

Contribution prompt A was applied to the complete definition and supplied source record, followed by the adversarial checks of prompt B. Local reproduction files are \texttt{batch025/checks.py}, \texttt{computations.json}, and the per-ID audit JSON. They were not public links at preparation time. Numerical checks are finite and do not prove unrestricted results. No LaTeX engine was found locally; compilation was not performed by this research agent.

\end{document}
